\section{Symmetries of the selected stationary germs}
\label{sec:stationary-germ-symmetry}
The Taylor construction of the output action uses the symmetries of
three finite-volume functions: the classical stationary value, the
completed determinant correction, and the inherited loop evaluated on
the stationary fine field. We prove these symmetries for the selected
stationary branch. In particular the determinant argument includes the
moving frame and its regularization.

Fix a block step, its spanning forest, a positive regulator $\mu$, and
the analytic stationary branch $U(V)$ in the original rooted tree
chart. The price gap is strictly positive. The standing hypotheses,
which are inputs to this section, are precisely the following: the
patch satisfies the branch compatibility of the kernel; the incoming
path satisfies its path bounds in the incoming clip; the branch
vanishes at zero, is analytic, and is stationary near the origin for
the tested action of \emph{every} choice of block representatives, not
only of one; the Hessian of every tested action in the tree-chart
variables at the origin has one positive lower bound; $\mu>0$; and a
local Haar chart with its density is fixed. Write $t$ for a choice of
block representatives. For a coarse field $V$ and a joint fine field
and frame $U$, let $S_t(V,U)$ be the classical action with the blocking
penalty. Let $H_t(V,U)$ be its Hessian in the joint left coordinates,
and let $F_t(U)$ be the completed gauge-fixing derivative. In the
notation of the determinant construction, put
\[
 M_t(V,U)=H_t(V,U)+F_t(U)^*F_t(U),\qquad
 \mathcal B_t(V,U)=\tfrac12\log\det M_t(V,U)
                                      -\log\det\Gamma(U).
\]
Here $\Gamma(U)=d_U^*d_U+C_\mu(U)$ is the same scalar gauge-orbit
operator, with $C_\mu=\mu^2P_\mu^9$. All matrices in this section
are real. The positive price gap and positive regulator make the
relevant determinants positive near the stationary origin, so these
expressions define analytic real germs there.

\subsection{Transport of the completed determinant}
A fine gauge transformation $k$ and coarse gauge transformation $u$
act on the joint field and frame. Their tangent maps in joint,
fine-link and scalar coordinates are orthogonal: they are products
of adjoint colour rotations and coordinate permutations. Denote the
joint tangent map by $T$ and the scalar map by $A$. Equivariance of
the block ribbon and of its frame logarithm gives
\[
                  F_t(U')T=AF_t(U).
\]
We verify why this identity includes both parts of the gauge fixing.
For its frame part, the logarithm near the identity obeys
$\log(khk^{-1})=\operatorname{Ad}_k\log h$; differentiating the
block-frame identity gives the intertwining of its derivative. For
its link part, the covariant difference intertwines the scalar and
link adjoint maps, and transposition gives the corresponding identity
for the adjoint difference. The scalar resolvent and hence its ninth
power intertwine because multiplication and inversion preserve an
intertwining relation. These are the two summands of $F_t$.

The classical action is gauge invariant. In left increment coordinates
the transformed action composed with $T$ equals the original action
near the increment origin. Twice differentiating this identity yields
\[
 T^*H_t(V',U')T=H_t(V,U).
\]
There is no derivative of $T$ in this calculation: the base fields
are fixed while differentiating the increment variable. Combining
this congruence with the gauge-fixing identity and $A^*A=I$ gives
\[
 T^*M_t(V',U')T=M_t(V,U).
\]
Similarly $\Gamma(U')$ is orthogonally conjugate to $\Gamma(U)$.
Taking determinants and using $(\det T)^2=(\det A)^2=1$ proves
\begin{equation}\label{eq:stationary-germ-gauge-det}
                   \mathcal B_t(V',U')=\mathcal B_t(V,U).
\end{equation}
The same transformation preserves $S_t$ and the inherited loop.

Signed lattice transformations require a slightly different tangent
map. For a transformed link $e$, let $e'$ be the original link it
reads. If its orientation is preserved, the tangent reads $X_{e'}$;
if reversed, it reads $-\operatorname{Ad}_{U_{e'}^{-1}}X_{e'}$.
This map is orthogonal and satisfies the exact increment identity
\[
 s^*(\exp(X)U)=\exp(T_UX)\,s^*U.
\]
The scalar tangent is the coordinate permutation; the joint tangent
also includes the corresponding frame permutation. The block
representative changes from $t$ to the transported choice $r$.
Equivariance of the ribbon, frame logarithm and covariant difference
therefore gives, with these tangent maps,
\[
 F_t(s^*U)T_U=AF_r(U),\qquad
 T_U^*H_t(s^*V,s^*U)T_U=H_r(V,U).
\]
For the penalty term this uses invariance of the group price under
inversion and conjugation: reversing both links in its relative
product leaves its price unchanged. The same resolvent argument
handles $C_\mu$. Taking determinants as above gives
\begin{equation}\label{eq:stationary-germ-signed-det}
                    \mathcal B_t(s^*V,s^*U)=\mathcal B_r(V,U).
\end{equation}
This is the completed moving determinant. It does not assert full
signed covariance of the separate auxiliary fine-lattice pole symbol.

\subsection{Returning a moved solution to the original chart}
A transformed stationary solution need not lie in the original tree
chart. Gauge it back using the tree gauge map. This map is analytic
near the trivial joint field, sends that field to itself, and leaves
the coarse input fixed when the coarse gauge transformation is the
identity. The transformed stationarity equation is still zero:
under the invertible tangent map its first derivative is the original
zero first derivative. The local uniqueness in the original chart,
proved by the stationary implicit-function construction, now identifies
the regauged solution with the selected solution at the transformed
coarse input. We use uniqueness only in that chart.

The classical value, the determinant correction by
\eqref{eq:stationary-germ-gauge-det}, and the inherited loop are
unchanged by this regauging. Consequently their values along any
analytic stationary probe through the trivial field equal their
values on the selected branch at that probe's coarse input. The
argument also applies to parameter-dependent gauge transformations:
freeze the parameter when comparing the action and matrices, and
then use analyticity and the local branch uniqueness on a common
neighborhood.

We must also remove the changed representative $r$ in
\eqref{eq:stationary-germ-signed-det}. The tested classical actions
for $t$ and $r$ agree as germs in both the coarse and tree-chart
variables, by the local branch construction of
Section~\ref{sec:stationary-action-prices}. Therefore their
stationary values and their tree-chart Hessians agree. The moving
Haar identity of Section~\ref{sec:ex-thermal-contact-cancellation}
expresses each thermal value as that tree Hessian determinant with
the same Haar factor, or equivalently as the same inherited loop
plus $\mathcal B_t$ or $\mathcal B_r$. Equality of the tree Hessians
and cancellation of the common inherited loop give equality of
$\mathcal B_t$ and $\mathcal B_r$ along the selected branch. Thus
both tested functions are independent of the representative as germs.
This uses equality on a neighborhood of the moving stationary point,
which follows from equality of the joint germs; equality at the
stationary value alone would not identify their Hessians.

\subsection{The identities used by the Taylor coefficients}
Apply the preceding argument first to a signed lattice transformation,
then to constant colour conjugation, and finally to a small variable
coarse gauge transformation. For the last operation the transformed
coarse logarithm is its actual nonlinear gauge action; transform the
joint frame accordingly and gauge the resulting stationary solution
back to the original chart. Each of the three analytic stationary
functions is unchanged as a germ under all these operations. The
inherited loop has the same invariances directly from its gauge
invariant atoms and finite signed average.

In summary, for every choice $t$ of block representatives each of the
three functions of the coarse link field, the classical stationary
value $E_t$, the completed determinant correction $\mathcal B_t$ along
the branch, and the inherited loop along the branch, is real analytic
at zero and is unchanged, as a germ at zero, by every signed lattice
transformation, by every constant colour conjugation, and by the
nonlinear action of small coarse gauge transformations. Moreover, for
two choices $t,r$ of representatives the germs of $E_t$ and $E_r$
coincide, and so do those of $\mathcal B_t$ and $\mathcal B_r$ along
the branch. All statements are for one forest. The comparison of the
stationary data of two different forests is a separate statement,
which is not proved in this section.

Differentiating these germ identities supplies the signed and colour
symmetries and the nonlinear gauge identities of the finite-volume
Taylor coefficients in
Section~\ref{sec:stationary-output-construction}. When the lower
positive-degree coefficients vanish, the differentiated nonlinear
gauge identity reduces to its leading linear Ward identity, as proved
there. All three functions here use the same stationary branch,
regulator and finite torus; no symmetry claim is imported from a
separately chosen stationary solution.
